%%% MA_A2_International_Law.tex -- reorganised; frames moved from the old MA_xx files (see comments)

%%% [moved from MA_2_International_Order.tex, frame 10]
\begin{frame}{Freedom of the Seas}
    \mytext{
    \begin{itemize}
        \item mare liberum (Hugo Grotius)
        \item 17th century
        \item British concept, 3 nautical miles (5.6 km)
        \item later de-facto 12 nautical miles
        \item ships under jurisdiction of flag they fly
    \end{itemize}
    }
\end{frame}

%%% [moved from MA_2_International_Order.tex, frame 11]
% \begin{frame}{The Gambia}
%     \mycomment{
%         To extend a rule from land to the sea just as far as you can look and fire a canon was sometimes even used the other way round. 
        
%         Look at the shape of The Gambia. It extends six miles left and right of the river Gambia---just as far as a British war ship was able to fire at the time. 
        
%         The Gambia is totally surrounded by Senegal, then under French rule. Still today, they use English and French as official languages, and have legal systems more leaning towards French written or English common law.
%     }
%     \showmybitmap[https://de.wikipedia.org/wiki/Datei:Gambia_rel_location_map.svg]{shared/bitmaps/Gambia_rel_location_map.svg.png}
% \end{frame}

%%% [moved from MA_2_International_Order.tex, frame 12]
\begin{frame}{UNCLOS}
    \mytext{
    \begin{itemize}
        \item Four treaties, 1958/1962
        \item 63 states, NATO, Soviet-bloc, OPEC etc.
        \item today 169 states, but not US
        \item Convention on the Territorial Sea and Continguous Zone
        \item Convention on the Continental Shelf
        \item Convention on the High Seas
        \item Convention on Fishing and Convention of Living Resources of the High Seas
        \item United Nations Convention on the Law of the Sea
    \end{itemize}
    }
\end{frame}

%%% [moved from MA_2_International_Order.tex, frame 13]
% \begin{frame}{International Waters and Exclusive Zones}
%     \showmytikz{../shared/tikz/tikz-waters.tex}
%     \mycomment{
%     There is a difference between exclusive right to economic use (of the waters, e.g. for fishing, or the resources at the sea bed) and legislation in regards to ships and traffic.

%     The legislation on board of the ship depends on  the flag under which the ship sails. However, contracts may specify which legislation applies, e.g. for work contracts.

%     \begin{equation}
%         1 NM = 1.852km
%     \end{equation}
%     A nautical mile is approximately one minute of arc of latitude.
%     }
% \end{frame}

%%% [moved from MA_2_International_Order.tex, frame 14]
\begin{frame}{International Waters and Exclusive Zones}
    \showmytikz{../shared/tikz/tikz-waters-two.tex}
    \mycomment{
    The territorial waters---and the contiguous zone---are ruled by the state just like its territory.

    Beyond that lies the freedom of the seas, and vessels have free passage. However, the state reserves the right to the sea bed, that is to lay cables, pipelines, drilling and so forth. This extends to the end of the continental shelf or to 200NM, whichever is greater.

    The exact shape and extend of both EEZ, the end of the continental shelf, and the extended continental shelf can give reason to disputes and conflict. Tiny islands often come with the benefit of large undersea territories.
    }
\end{frame}

%%% [moved from MA_2_International_Order.tex, frame 19]
\begin{frame}{Security and Military Alliances}
    \mytext{
    \begin{itemize}
        \item OSCE Organisation for Security and Co-operation in Europe (Vienna)
        \begin{itemize}
            \item All states in Europe  
            \item Conference on security
        \end{itemize}
        \item NATO 
        \item Council of Europe
            \begin{itemize}
                \item Not part of EU system
                \item European Court of Human Rights
                \item European Convention on Human Rights
            \end{itemize}        
    \end{itemize}
    }
\end{frame}

